\documentclass{beamer}
\usetheme{Madrid}
\usecolortheme{default}
\usepackage[utf8]{inputenc}
\usepackage{amsmath}
\usepackage{amsfonts}
\usepackage{amssymb}
\usepackage{tikz}

\title{CRISC Domain A - 100 Impossible Hard Questions}
\subtitle{Governance - Expert Level That Will Break You}
\author{CRISC Mastery Series}
\date{\today}

\begin{document}
	
	\begin{frame}
		\titlepage
	\end{frame}
	
	\begin{frame}{Warning}
		\begin{alertblock}{Extreme Difficulty}
			These questions test the deepest governance judgment.\\
			Every option will look plausible.\\
			You must think like a board advisor, not a technician.
		\end{alertblock}
	\end{frame}
	
	\begin{frame}{Outline}
		\tableofcontents
	\end{frame}
	
	% ============================================================
	\section{1. Organizational Governance}
	% ============================================================
	
	\begin{frame}{Q1: Governance vs Management}
		\begin{block}{Question}
			The board sets a policy requiring security review for all new vendors. The CIO writes the procedures and assigns reviewers. Who is performing governance and who is performing management?
		\end{block}
		\begin{enumerate}[(A)]
			\item Both are governance
			\item Both are management
			\item Board performs governance; CIO performs management
			\item Board performs management; CIO performs governance
		\end{enumerate}
	\end{frame}
	
	\begin{frame}{Q1: Explanation}
		\begin{block}{Correct Answer: C}
			\textbf{Governance vs Management:}
			\begin{itemize}
				\item Governance sets direction and accountability
				\item Management executes within that direction
				\item Board policy is governance
				\item CIO procedures are management
			\end{itemize}
		\end{block}
		\begin{alertblock}{Critical Trap}
			Governance directs; management executes.
		\end{alertblock}
	\end{frame}
	
	\begin{frame}{Q2: Accountability Trap}
		\begin{block}{Question}
			The CISO delegates incident response to the SOC manager. The SOC manager mishandles a breach. Who is accountable to the board?
		\end{block}
		\begin{enumerate}[(A)]
			\item The SOC analyst who missed the alert
			\item The SOC manager
			\item The CISO
			\item The vendor who sold the SIEM
		\end{enumerate}
	\end{frame}
	
	\begin{frame}{Q2: Explanation}
		\begin{block}{Correct Answer: C}
			\textbf{Accountability does not delegate:}
			\begin{itemize}
				\item Delegation transfers tasks, not accountability
				\item The CISO remains accountable for the security program
				\item The SOC manager is responsible for operations
				\item Accountability flows upward
			\end{itemize}
		\end{block}
		\begin{alertblock}{Critical Trap}
			Accountability does not delegate with the task.
		\end{alertblock}
	\end{frame}
	
	\begin{frame}{Q3: Policy vs Standard}
		\begin{block}{Question}
			A statement says ``all customer data must be encrypted.'' A second document says ``AES-256 with FIPS 140-2 modules.'' Which is the policy and which is the standard?
		\end{block}
		\begin{enumerate}[(A)]
			\item Both are policies
			\item Both are standards
			\item The first is the policy; the second is the standard
			\item The first is the standard; the second is the policy
		\end{enumerate}
	\end{frame}
	
	\begin{frame}{Q3: Explanation}
		\begin{block}{Correct Answer: C}
			\textbf{Policy vs Standard:}
			\begin{itemize}
				\item Policy states intent at a high level
				\item Standards are mandatory, specific requirements
				\item The encryption requirement is policy
				\item The algorithm specification is standard
			\end{itemize}
		\end{block}
		\begin{alertblock}{Critical Trap}
			Policy states intent; standards specify requirements.
		\end{alertblock}
	\end{frame}
	
	\begin{frame}{Q4: Tone at the Top}
		\begin{block}{Question}
			Senior executives publicly commit to risk management but cut the security budget each year. What does this signal to the organization?
		\end{block}
		\begin{enumerate}[(A)]
			\item Strong risk culture
			\item The culture will mirror actions, not words
			\item Budget cuts are always appropriate
			\item The board should relax oversight
		\end{enumerate}
	\end{frame}
	
	\begin{frame}{Q4: Explanation}
		\begin{block}{Correct Answer: B}
			\textbf{Actions shape culture:}
			\begin{itemize}
				\item Actions shape culture more than statements
				\item Employees observe what leaders do
				\item Inconsistent messages erode trust
			\end{itemize}
		\end{block}
		\begin{alertblock}{Critical Trap}
			Actions speak louder than commitments.
		\end{alertblock}
	\end{frame}
	
	\begin{frame}{Q5: Risk Appetite Authority}
		\begin{block}{Question}
			The risk manager proposes the risk appetite. Who approves it?
		\end{block}
		\begin{enumerate}[(A)]
			\item The risk manager
			\item The audit committee alone
			\item The board of directors
			\item The IT steering committee
		\end{enumerate}
	\end{frame}
	
	\begin{frame}{Q5: Explanation}
		\begin{block}{Correct Answer: C}
			\textbf{Appetite is board-owned:}
			\begin{itemize}
				\item Appetite is a board-level decision
				\item It reflects strategy and stakeholder interests
				\item Management proposes; the board approves
			\end{itemize}
		\end{block}
		\begin{alertblock}{Critical Trap}
			Appetite is board-owned.
		\end{alertblock}
	\end{frame}
	
	\begin{frame}{Q6: Three Lines - First}
		\begin{block}{Question}
			Which function is the FIRST line of defense?
		\end{block}
		\begin{enumerate}[(A)]
			\item Risk management function
			\item Internal audit
			\item Business operations
			\item External audit
		\end{enumerate}
	\end{frame}
	
	\begin{frame}{Q6: Explanation}
		\begin{block}{Correct Answer: C}
			\textbf{First line = business operations:}
			\begin{itemize}
				\item The first line owns and manages risk
				\item The second line oversees and advises
				\item The third line provides independent assurance
			\end{itemize}
		\end{block}
		\begin{alertblock}{Critical Trap}
			The first line owns risk.
		\end{alertblock}
	\end{frame}
	
	\begin{frame}{Q7: Three Lines - Second}
		\begin{block}{Question}
			Who is the SECOND line of defense?
		\end{block}
		\begin{enumerate}[(A)]
			\item Operational managers
			\item External auditors
			\item A risk and compliance function that oversees first-line activities
			\item The board
		\end{enumerate}
	\end{frame}
	
	\begin{frame}{Q7: Explanation}
		\begin{block}{Correct Answer: C}
			\textbf{Second line oversees:}
			\begin{itemize}
				\item The second line oversees risk practices
				\item It advises and challenges the first line
				\item The third line provides independent assurance
			\end{itemize}
		\end{block}
		\begin{alertblock}{Critical Trap}
			Second line oversees, third line audits.
		\end{alertblock}
	\end{frame}
	
	\begin{frame}{Q8: Three Lines - Third}
		\begin{block}{Question}
			What is the MAIN purpose of the third line of defense?
		\end{block}
		\begin{enumerate}[(A)]
			\item Operating controls
			\item Setting risk appetite
			\item Independent assurance over the first and second lines
			\item Approving risk responses
		\end{enumerate}
	\end{frame}
	
	\begin{frame}{Q8: Explanation}
		\begin{block}{Correct Answer: C}
			\textbf{Third line = independent assurance:}
			\begin{itemize}
				\item Third line is independent assurance
				\item It does not own risk
				\item It evaluates effectiveness
			\end{itemize}
		\end{block}
		\begin{alertblock}{Critical Trap}
			Third line provides independent assurance.
		\end{alertblock}
	\end{frame}
	
	\begin{frame}{Q9: Board Oversight}
		\begin{block}{Question}
			A material IT risk exceeds appetite and is not being treated. What is the board's BEST action?
		\end{block}
		\begin{enumerate}[(A)]
			\item Implement the control directly
			\item Direct management to treat or formally accept the risk with documentation
			\item Ignore it
			\item Delegate to the vendor
		\end{enumerate}
	\end{frame}
	
	\begin{frame}{Q9: Explanation}
		\begin{block}{Correct Answer: B}
			\textbf{Boards direct; management executes:}
			\begin{itemize}
				\item Boards direct management
				\item They do not implement controls
				\item Formal treatment or acceptance is required
			\end{itemize}
		\end{block}
		\begin{alertblock}{Critical Trap}
			Boards direct; management executes.
		\end{alertblock}
	\end{frame}
	
	\begin{frame}{Q10: Risk Committee}
		\begin{block}{Question}
			An IT risk committee meets quarterly but decisions never reach the board. What is the MAIN governance gap?
		\end{block}
		\begin{enumerate}[(A)]
			\item Too many meetings
			\item Insufficient reporting escalation to the board
			\item Too many committee members
			\item Not enough coffee
		\end{enumerate}
	\end{frame}
	
	\begin{frame}{Q10: Explanation}
		\begin{block}{Correct Answer: B}
			\textbf{Escalation gap:}
			\begin{itemize}
				\item Decisions must reach the accountable body
				\item Escalation keeps governance connected
				\item Committee work without board visibility is ineffective
			\end{itemize}
		\end{block}
		\begin{alertblock}{Critical Trap}
			Without escalation, governance is disconnected.
		\end{alertblock}
	\end{frame}
	
	\begin{frame}{Q11: RACI Accountable}
		\begin{block}{Question}
			In a RACI chart, who is the single accountable party for an activity?
		\end{block}
		\begin{enumerate}[(A)]
			\item Everyone in the ``R'' column
			\item The ``A'' - one accountable person
			\item The ``C'' - consulted party
			\item The ``I'' - informed party
		\end{enumerate}
	\end{frame}
	
	\begin{frame}{Q11: Explanation}
		\begin{block}{Correct Answer: B}
			\textbf{One accountable per activity:}
			\begin{itemize}
				\item ``A'' is single-point accountability
				\item ``R'' does the work
				\item ``C'' advises; ``I'' is informed
			\end{itemize}
		\end{block}
		\begin{alertblock}{Critical Trap}
			One accountable, one ``A''.
		\end{alertblock}
	\end{frame}
	
	\begin{frame}{Q12: RACI Consulted}
		\begin{block}{Question}
			When an architecture team is marked ``Consulted'' on a risk initiative, what does that mean?
		\end{block}
		\begin{enumerate}[(A)]
			\item They perform the work
			\item They approve the decision
			\item They provide subject matter expertise before decisions
			\item They receive updates only
		\end{enumerate}
	\end{frame}
	
	\begin{frame}{Q12: Explanation}
		\begin{block}{Correct Answer: C}
			\textbf{Consulted = advice:}
			\begin{itemize}
				\item Consulted parties give input
				\item Responsible parties do the work
				\item Accountable parties own decisions
			\end{itemize}
		\end{block}
		\begin{alertblock}{Critical Trap}
			Consulted means advice, not decision.
		\end{alertblock}
	\end{frame}
	
	\begin{frame}{Q13: Enterprise Risk Committee}
		\begin{block}{Question}
			The enterprise risk committee includes only IT leaders. What is the MOST likely weakness?
		\end{block}
		\begin{enumerate}[(A)]
			\item Too many technology voices, not enough business representation
			\item Too many business voices
			\item Too few vendors
			\item Too many meetings
		\end{enumerate}
	\end{frame}
	
	\begin{frame}{Q13: Explanation}
		\begin{block}{Correct Answer: A}
			\textbf{Cross-functional needed:}
			\begin{itemize}
				\item Enterprise risk requires cross-functional input
				\item IT alone cannot represent business risk
				\item Business, legal, finance, and operations should be included
			\end{itemize}
		\end{block}
		\begin{alertblock}{Critical Trap}
			Enterprise risk needs cross-functional voices.
		\end{alertblock}
	\end{frame}
	
	\begin{frame}{Q14: Governance of AI}
		\begin{block}{Question}
			The organization adopts an AI tool without a defined owner or accountability. What governance principle is violated?
		\end{block}
		\begin{enumerate}[(A)]
			\item Segregation of duties
			\item Accountability must be assigned
			\item Least privilege
			\item Defense in depth
		\end{enumerate}
	\end{frame}
	
	\begin{frame}{Q14: Explanation}
		\begin{block}{Correct Answer: B}
			\textbf{Accountability required:}
			\begin{itemize}
				\item Every initiative needs an accountable owner
				\item Unowned technology creates unmanaged risk
				\item Governance requires explicit accountability
			\end{itemize}
		\end{block}
		\begin{alertblock}{Critical Trap}
			No owner means no governance.
		\end{alertblock}
	\end{frame}
	
	\begin{frame}{Q15: Policy Exceptions}
		\begin{block}{Question}
			An executive grants a policy exception verbally to a peer. What is the governance flaw?
		\end{block}
		\begin{enumerate}[(A)]
			\item No flaw, verbal is fine
			\item Exceptions must be documented, approved by the appropriate authority, and reviewed
			\item Exceptions should be free
			\item Policy is optional
		\end{enumerate}
	\end{frame}
	
	\begin{frame}{Q15: Explanation}
		\begin{block}{Correct Answer: B}
			\textbf{Exceptions need documentation:}
			\begin{itemize}
				\item Exceptions need documentation
				\item Approval authority must be defined
				\item Review ensures validity
			\end{itemize}
		\end{block}
		\begin{alertblock}{Critical Trap}
			Verbal exceptions are ungovernable.
		\end{alertblock}
	\end{frame}
	
	% ============================================================
	\section{2. Legal, Regulatory and Ethical Requirements}
	% ============================================================
	
	\begin{frame}{Q16: Regulatory Prioritization}
		\begin{block}{Question}
			The organization operates in multiple jurisdictions with conflicting privacy laws. What should the risk practitioner recommend?
		\end{block}
		\begin{enumerate}[(A)]
			\item Follow one jurisdiction only
			\item Identify the intersection of obligations and design controls that meet the strictest applicable requirements
			\item Ignore the conflict
			\item Choose based on cost
		\end{enumerate}
	\end{frame}
	
	\begin{frame}{Q16: Explanation}
		\begin{block}{Correct Answer: B}
			\textbf{Compliance is a floor:}
			\begin{itemize}
				\item Conflicting laws require a control floor
				\item The strictest requirement sets the baseline
				\item Legal review confirms applicability
			\end{itemize}
		\end{block}
		\begin{alertblock}{Critical Trap}
			Compliance is a floor, not a menu.
		\end{alertblock}
	\end{frame}
	
	\begin{frame}{Q17: Ethics Beyond Law}
		\begin{block}{Question}
			An activity is legal but violates the ISACA Code of Ethics. What should the practitioner do?
		\end{block}
		\begin{enumerate}[(A)]
			\item Proceed because it is legal
			\item Refuse and escalate, citing professional ethics
			\item Ignore the code
			\item Ask a peer
		\end{enumerate}
	\end{frame}
	
	\begin{frame}{Q17: Explanation}
		\begin{block}{Correct Answer: B}
			\textbf{Ethics is independent of legality:}
			\begin{itemize}
				\item Ethics is independent of legality
				\item Certification holders commit to the Code
				\item Escalation preserves professional standards
			\end{itemize}
		\end{block}
		\begin{alertblock}{Critical Trap}
			Legal is not the same as ethical.
		\end{alertblock}
	\end{frame}
	
	\begin{frame}{Q18: Whistleblower Protection}
		\begin{block}{Question}
			An employee reports fraud but is retaliated against. What governance failure is MOST visible?
		\end{block}
		\begin{enumerate}[(A)]
			\item Weak anti-retaliation controls and culture
			\item Weak marketing controls
			\item Weak IT procurement
			\item Weak training budget
		\end{enumerate}
	\end{frame}
	
	\begin{frame}{Q18: Explanation}
		\begin{block}{Correct Answer: A}
			\textbf{Retaliation kills reporting:}
			\begin{itemize}
				\item Retaliation kills honest escalation
				\item Anonymous channels and protections are required
				\item Culture signals the absence of safety
			\end{itemize}
		\end{block}
		\begin{alertblock}{Critical Trap}
			Silence is a governance failure.
		\end{alertblock}
	\end{frame}
	
	\begin{frame}{Q19: Regulatory Ambiguity}
		\begin{block}{Question}
			A regulator has issued guidance but not a formal rule. How should the practitioner respond?
		\end{block}
		\begin{enumerate}[(A)]
			\item Wait for a formal rule
			\item Assess exposure, interpret guidance with legal, and prepare proportionate controls
			\item Ignore guidance
			\item Move to another jurisdiction
		\end{enumerate}
	\end{frame}
	
	\begin{frame}{Q19: Explanation}
		\begin{block}{Correct Answer: B}
			\textbf{Guidance signals intent:}
			\begin{itemize}
				\item Guidance signals regulatory intent
				\item Legal counsel interprets applicability
				\item Preparation reduces future risk
			\end{itemize}
		\end{block}
		\begin{alertblock}{Critical Trap}
			Treat guidance as a signal, not silence.
		\end{alertblock}
	\end{frame}
	
	\begin{frame}{Q20: Ethics of Automation}
		\begin{block}{Question}
			An AI model automates decisions with material impact. What should governance require before deployment?
		\end{block}
		\begin{enumerate}[(A)]
			\item Faster deployment
			\item Explainability, oversight, and a defined accountable owner
			\item Fewer controls
			\item No user notification
		\end{enumerate}
	\end{frame}
	
	\begin{frame}{Q20: Explanation}
		\begin{block}{Correct Answer: B}
			\textbf{Accountable automation:}
			\begin{itemize}
				\item Explainability supports accountability
				\item Human oversight catches bias and failure
				\item Unowned models cannot be governed
			\end{itemize}
		\end{block}
		\begin{alertblock}{Critical Trap}
			Automation requires accountability, not just speed.
		\end{alertblock}
	\end{frame}
	
	\begin{frame}{Q21: Third-Party Ethics}
		\begin{block}{Question}
			A vendor follows local laws but not the organization's ethical standards. What should governance require?
		\end{block}
		\begin{enumerate}[(A)]
			\item Accept local law as the standard
			\item Apply the organization's ethical standards contractually or reassess the relationship
			\item Ignore the difference
			\item Outsource more
		\end{enumerate}
	\end{frame}
	
	\begin{frame}{Q21: Explanation}
		\begin{block}{Correct Answer: B}
			\textbf{Ethics follow the enterprise:}
			\begin{itemize}
				\item Ethics are organizational, not jurisdictional
				\item Contracts can require compliance with internal standards
				\item Mismatch is a governance gap
			\end{itemize}
		\end{block}
		\begin{alertblock}{Critical Trap}
			Ethical standards follow the enterprise.
		\end{alertblock}
	\end{frame}
	
	\begin{frame}{Q22: Data Subject Rights}
		\begin{block}{Question}
			A data subject requests deletion of their data. What governance principle MOST applies?
		\end{block}
		\begin{enumerate}[(A)]
			\item Availability first
			\item Data subject rights must be honored per regulation and policy
			\item Profit first
			\item Delay as long as possible
		\end{enumerate}
	\end{frame}
	
	\begin{frame}{Q22: Explanation}
		\begin{block}{Correct Answer: B}
			\textbf{Rights are obligations:}
			\begin{itemize}
				\item Data rights are a regulatory obligation
				\item Policies must define handling
				\item Delay creates legal exposure
			\end{itemize}
		\end{block}
		\begin{alertblock}{Critical Trap}
			Rights are obligations, not options.
		\end{alertblock}
	\end{frame}
	
	\begin{frame}{Q23: Contractual Obligations}
		\begin{block}{Question}
			A contract requires 24-hour breach notification, but policy says 72 hours. What should the practitioner do?
		\end{block}
		\begin{enumerate}[(A)]
			\item Follow policy
			\item Identify the gap and recommend amending policy or contract to align
			\item Ignore the contract
			\item Delay notification
		\end{enumerate}
	\end{frame}
	
	\begin{frame}{Q23: Explanation}
		\begin{block}{Correct Answer: B}
			\textbf{Contracts bind:}
			\begin{itemize}
				\item Contracts are binding obligations
				\item Conflicts must be resolved in governance
				\item Notification timing matters legally
			\end{itemize}
		\end{block}
		\begin{alertblock}{Critical Trap}
			Contractual timing overrides internal defaults.
		\end{alertblock}
	\end{frame}
	
	\begin{frame}{Q24: Legal Counsel Role}
		\begin{block}{Question}
			When should the risk practitioner involve legal counsel?
		\end{block}
		\begin{enumerate}[(A)]
			\item Only after an incident
			\item When regulatory interpretation, contractual risk, or legal exposure is unclear
			\item Never
			\item Only at year-end
		\end{enumerate}
	\end{frame}
	
	\begin{frame}{Q24: Explanation}
		\begin{block}{Correct Answer: B}
			\textbf{Engage legal early:}
			\begin{itemize}
				\item Legal interprets obligations
				\item Timely involvement avoids missteps
				\item Year-end reviews are too late
			\end{itemize}
		\end{block}
		\begin{alertblock}{Critical Trap}
			Engage legal early on legal questions.
		\end{alertblock}
	\end{frame}
	
	\begin{frame}{Q25: Ethics Reporting}
		\begin{block}{Question}
			The organization lacks an anonymous ethics reporting channel. What is the MAIN governance gap?
		\end{block}
		\begin{enumerate}[(A)]
			\item Too few channels for reporting wrongdoing safely
			\item Too many meetings
			\item Too much marketing
			\item Too many tools
		\end{enumerate}
	\end{frame}
	
	\begin{frame}{Q25: Explanation}
		\begin{block}{Correct Answer: A}
			\textbf{Safe channels needed:}
			\begin{itemize}
				\item Safe reporting supports transparency
				\item Anonymous channels protect reporters
				\item Absence increases undetected issues
			\end{itemize}
		\end{block}
		\begin{alertblock}{Critical Trap}
			Safe reporting is a governance control.
		\end{alertblock}
	\end{frame}
	
	% ============================================================
	\section{3. Roles and Responsibilities}
	% ============================================================
	
	\begin{frame}{Q26: CRO Role}
		\begin{block}{Question}
			Which is the PRIMARY responsibility of the Chief Risk Officer (CRO)?
		\end{block}
		\begin{enumerate}[(A)]
			\item Operating IT controls
			\item Owning enterprise risk strategy and oversight
			\item Auditing the enterprise
			\item Approving budgets
		\end{enumerate}
	\end{frame}
	
	\begin{frame}{Q26: Explanation}
		\begin{block}{Correct Answer: B}
			\textbf{CRO sets strategy:}
			\begin{itemize}
				\item The CRO leads enterprise risk strategy
				\item The CRO does not operate controls
				\item The CRO advises and oversees
			\end{itemize}
		\end{block}
		\begin{alertblock}{Critical Trap}
			The CRO sets strategy and oversees.
		\end{alertblock}
	\end{frame}
	
	\begin{frame}{Q27: CISO Role}
		\begin{block}{Question}
			A CISO reports to the CIO who also owns IT operations. What potential conflict arises?
		\end{block}
		\begin{enumerate}[(A)]
			\item No conflict
			\item Independence is weakened; the CISO may be less able to challenge IT operations
			\item Stronger security
			\item Better governance
		\end{enumerate}
	\end{frame}
	
	\begin{frame}{Q27: Explanation}
		\begin{block}{Correct Answer: B}
			\textbf{Independence matters:}
			\begin{itemize}
				\item Independence is critical for security oversight
				\item Reporting to operations owners can mute challenges
				\item Some organizations place the CISO under the CEO or CRO
			\end{itemize}
		\end{block}
		\begin{alertblock}{Critical Trap}
			Independence shapes honest challenge.
		\end{alertblock}
	\end{frame}
	
	\begin{frame}{Q28: Data Owner Role}
		\begin{block}{Question}
			Who is accountable for determining classification and access to customer data?
		\end{block}
		\begin{enumerate}[(A)]
			\item Data custodian
			\item Data owner
			\item Data user
			\item System administrator
		\end{enumerate}
	\end{frame}
	
	\begin{frame}{Q28: Explanation}
		\begin{block}{Correct Answer: B}
			\textbf{Owners classify:}
			\begin{itemize}
				\item Data owners classify and authorize access
				\item Custodians implement controls
				\item Administrators operate systems
			\end{itemize}
		\end{block}
		\begin{alertblock}{Critical Trap}
			Owners classify and authorize; custodians implement.
		\end{alertblock}
	\end{frame}
	
	\begin{frame}{Q29: Data Custodian Role}
		\begin{block}{Question}
			Who implements and maintains the technical controls protecting data?
		\end{block}
		\begin{enumerate}[(A)]
			\item Data owner
			\item Data custodian
			\item The board
			\item The regulator
		\end{enumerate}
	\end{frame}
	
	\begin{frame}{Q29: Explanation}
		\begin{block}{Correct Answer: B}
			\textbf{Custodians implement:}
			\begin{itemize}
				\item Custodians implement controls
				\item Owners define requirements
				\item Boards oversee, not implement
			\end{itemize}
		\end{block}
		\begin{alertblock}{Critical Trap}
			Custodians implement; owners decide.
		\end{alertblock}
	\end{frame}
	
	\begin{frame}{Q30: Business Unit Accountability}
		\begin{block}{Question}
			A business unit launches a new service without security review. Who is accountable?
		\end{block}
		\begin{enumerate}[(A)]
			\item The vendor
			\item The business unit leader who owns the service
			\item The IT helpdesk
			\item The external auditor
		\end{enumerate}
	\end{frame}
	
	\begin{frame}{Q30: Explanation}
		\begin{block}{Correct Answer: B}
			\textbf{Business leaders own business risk:}
			\begin{itemize}
				\item Business leaders own business risk
				\item IT and security advise
				\item The vendor is not accountable to the enterprise
			\end{itemize}
		\end{block}
		\begin{alertblock}{Critical Trap}
			Business leaders own business risk.
		\end{alertblock}
	\end{frame}
	
	\begin{frame}{Q31: Steering Committee}
		\begin{block}{Question}
			What is the PRIMARY purpose of an IT steering committee?
		\end{block}
		\begin{enumerate}[(A)]
			\item Writing code
			\item Aligning IT investments and priorities with business objectives
			\item Auditing IT
			\item Operating infrastructure
		\end{enumerate}
	\end{frame}
	
	\begin{frame}{Q31: Explanation}
		\begin{block}{Correct Answer: B}
			\textbf{Alignment purpose:}
			\begin{itemize}
				\item Steering committees govern IT priorities
				\item Alignment ensures value
				\item Operational work belongs elsewhere
			\end{itemize}
		\end{block}
		\begin{alertblock}{Critical Trap}
			Steering committees align IT and business.
		\end{alertblock}
	\end{frame}
	
	\begin{frame}{Q32: Board Audit Committee}
		\begin{block}{Question}
			Which activity BEST fits the board audit committee?
		\end{block}
		\begin{enumerate}[(A)]
			\item Implementing controls
			\item Overseeing financial reporting, internal control, and risk assurance
			\item Operating the IT service desk
			\item Selecting vendors
		\end{enumerate}
	\end{frame}
	
	\begin{frame}{Q32: Explanation}
		\begin{block}{Correct Answer: B}
			\textbf{Audit committees oversee:}
			\begin{itemize}
				\item Audit committees oversee assurance
				\item They do not implement or operate
				\item Independence is essential
			\end{itemize}
		\end{block}
		\begin{alertblock}{Critical Trap}
			Audit committees oversee, not operate.
		\end{alertblock}
	\end{frame}
	
	\begin{frame}{Q33: Delegation Limits}
		\begin{block}{Question}
			Can a board permanently delegate risk appetite to a committee?
		\end{block}
		\begin{enumerate}[(A)]
			\item Yes, fully
			\item No; the board retains ultimate accountability
			\item Yes, if the CEO agrees
			\item Yes, if regulations allow
		\end{enumerate}
	\end{frame}
	
	\begin{frame}{Q33: Explanation}
		\begin{block}{Correct Answer: B}
			\textbf{Ultimate accountability stays:}
			\begin{itemize}
				\item Delegation can be operational
				\item Ultimate accountability remains with the board
				\item Appetite is board-level
			\end{itemize}
		\end{block}
		\begin{alertblock}{Critical Trap}
			Ultimate accountability stays with the board.
		\end{alertblock}
	\end{frame}
	
	\begin{frame}{Q34: Segregation in Governance}
		\begin{block}{Question}
			The same executive approves risk and audits performance. What is the concern?
		\end{block}
		\begin{enumerate}[(A)]
			\item Efficiency
			\item Lack of independence and potential conflict of interest
			\item Cost savings
			\item Better documentation
		\end{enumerate}
	\end{frame}
	
	\begin{frame}{Q34: Explanation}
		\begin{block}{Correct Answer: B}
			\textbf{Independence prevents self-review:}
			\begin{itemize}
				\item Independence prevents self-review
				\item Different parties should audit and approve
				\item Combining roles weakens governance
			\end{itemize}
		\end{block}
		\begin{alertblock}{Critical Trap}
			Self-review weakens governance.
		\end{alertblock}
	\end{frame}
	
	\begin{frame}{Q35: Governance of Third Parties}
		\begin{block}{Question}
			Who should approve critical vendor onboarding?
		\end{block}
		\begin{enumerate}[(A)]
			\item The procurement clerk
			\item The business owner, supported by risk and security review
			\item Any team member
			\item The vendor itself
		\end{enumerate}
	\end{frame}
	
	\begin{frame}{Q35: Explanation}
		\begin{block}{Correct Answer: B}
			\textbf{Owners approve:}
			\begin{itemize}
				\item Business owners accept the risk
				\item Risk and security advise
				\item Approval needs authority and context
			\end{itemize}
		\end{block}
		\begin{alertblock}{Critical Trap}
			Owners approve; advisors advise.
		\end{alertblock}
	\end{frame}
	
	% ============================================================
	\section{4. Enterprise Risk Management}
	% ============================================================
	
	\begin{frame}{Q36: ERM Scope}
		\begin{block}{Question}
			Enterprise risk management should cover:
		\end{block}
		\begin{enumerate}[(A)]
			\item Only IT risks
			\item Only financial risks
			\item All material risks across the enterprise
			\item Only compliance risks
		\end{enumerate}
	\end{frame}
	
	\begin{frame}{Q36: Explanation}
		\begin{block}{Correct Answer: C}
			\textbf{ERM is enterprise-wide:}
			\begin{itemize}
				\item ERM is enterprise-wide
				\item It aggregates and prioritizes risks
				\item Siloed risk management misses cross-cutting exposure
			\end{itemize}
		\end{block}
		\begin{alertblock}{Critical Trap}
			ERM is enterprise-wide.
		\end{alertblock}
	\end{frame}
	
	\begin{frame}{Q37: Risk Universe}
		\begin{block}{Question}
			What is the ``risk universe''?
		\end{block}
		\begin{enumerate}[(A)]
			\item All identified IT assets
			\item The full set of material risks facing the enterprise
			\item A list of vendors
			\item A list of projects
		\end{enumerate}
	\end{frame}
	
	\begin{frame}{Q37: Explanation}
		\begin{block}{Correct Answer: B}
			\textbf{The risk universe is enterprise-wide:}
			\begin{itemize}
				\item The risk universe is enterprise-wide
				\item It includes strategic, operational, financial, compliance, and IT risks
				\item It is the basis for prioritization
			\end{itemize}
		\end{block}
		\begin{alertblock}{Critical Trap}
			The universe is the full set of material risks.
		\end{alertblock}
	\end{frame}
	
	\begin{frame}{Q38: Risk Portfolio}
		\begin{block}{Question}
			Risk portfolio thinking means:
		\end{block}
		\begin{enumerate}[(A)]
			\item Assessing risks only individually
			\item Considering the aggregate and interactions of risks together
			\item Ignoring correlations
			\item Ranking risks by age
		\end{enumerate}
	\end{frame}
	
	\begin{frame}{Q38: Explanation}
		\begin{block}{Correct Answer: B}
			\textbf{Portfolio view:}
			\begin{itemize}
				\item Portfolios capture combined exposure
				\item Correlations and dependencies matter
				\item Single-risk views miss systemic risk
			\end{itemize}
		\end{block}
		\begin{alertblock}{Critical Trap}
			Portfolio view catches combined exposure.
		\end{alertblock}
	\end{frame}
	
	\begin{frame}{Q39: Risk Aggregation}
		\begin{block}{Question}
			Why is enterprise-level risk aggregation important?
		\end{block}
		\begin{enumerate}[(A)]
			\item It hides small risks
			\item It reveals combined exposure that single-risk views may miss
			\item It removes the need for controls
			\item It replaces business judgment
		\end{enumerate}
	\end{frame}
	
	\begin{frame}{Q39: Explanation}
		\begin{block}{Correct Answer: B}
			\textbf{Aggregation shows total exposure:}
			\begin{itemize}
				\item Aggregation shows total exposure
				\item Multiple small risks can exceed appetite
				\item It supports informed strategy
			\end{itemize}
		\end{block}
		\begin{alertblock}{Critical Trap}
			Small risks can add up to big exposure.
		\end{alertblock}
	\end{frame}
	
	\begin{frame}{Q40: Risk Tolerance Statements}
		\begin{block}{Question}
			A tolerance statement says ``risk above \$5M requires board approval.'' What is this an example of?
		\end{block}
		\begin{enumerate}[(A)]
			\item A control
			\item A defined escalation threshold reflecting tolerance
			\item An audit finding
			\item A project plan
		\end{enumerate}
	\end{frame}
	
	\begin{frame}{Q40: Explanation}
		\begin{block}{Correct Answer: B}
			\textbf{Tolerance expressed as threshold:}
			\begin{itemize}
				\item Tolerance is expressed in thresholds
				\item Thresholds drive escalation
				\item Clarity improves governance
			\end{itemize}
		\end{block}
		\begin{alertblock}{Critical Trap}
			Tolerance is expressed in clear thresholds.
		\end{alertblock}
	\end{frame}
	
	\begin{frame}{Q41: Risk Reporting Cadence}
		\begin{block}{Question}
			Enterprise risk reporting to the board should be:
		\end{block}
		\begin{enumerate}[(A)]
			\item Only during audits
			\item On a defined cadence with escalation for material changes
			\item Never
			\item Annually only
		\end{enumerate}
	\end{frame}
	
	\begin{frame}{Q41: Explanation}
		\begin{block}{Correct Answer: B}
			\textbf{Cadence plus escalation:}
			\begin{itemize}
				\item Cadence supports oversight
				\item Escalation handles urgent matters
				\item Annual alone misses changes
			\end{itemize}
		\end{block}
		\begin{alertblock}{Critical Trap}
			Steady cadence plus escalation.
		\end{alertblock}
	\end{frame}
	
	\begin{frame}{Q42: Risk Culture Measurement}
		\begin{block}{Question}
			Which is a good signal that risk culture is healthy?
		\end{block}
		\begin{enumerate}[(A)]
			\item Employees fear reporting risks
			\item Employees report risks early without fear
			\item No risks are ever reported
			\item Only executives report
		\end{enumerate}
	\end{frame}
	
	\begin{frame}{Q42: Explanation}
		\begin{block}{Correct Answer: B}
			\textbf{Early reporting = healthy culture:}
			\begin{itemize}
				\item Early reporting shows trust and awareness
				\item Fear suppresses information
				\item Silence signals cultural failure
			\end{itemize}
		\end{block}
		\begin{alertblock}{Critical Trap}
			Early reporting = healthy culture.
		\end{alertblock}
	\end{frame}
	
	\begin{frame}{Q43: Risk Appetite Translation}
		\begin{block}{Question}
			A board appetite statement says ``avoid regulatory penalties.'' How should it translate into IT risk thresholds?
		\end{block}
		\begin{enumerate}[(A)]
			\item No thresholds needed
			\item Define specific compliance KRIs and thresholds that reveal regulatory risk
			\item Ignore regulatory risk
			\item Rely on the legal team
		\end{enumerate}
	\end{frame}
	
	\begin{frame}{Q43: Explanation}
		\begin{block}{Correct Answer: B}
			\textbf{Translate intent:}
			\begin{itemize}
				\item Appetite must translate to measurable indicators
				\item KRIs turn intent into oversight
				\item Reliance on one function without thresholds leaves gaps
			\end{itemize}
		\end{block}
		\begin{alertblock}{Critical Trap}
			Translate intent into measurable indicators.
		\end{alertblock}
	\end{frame}
	
	\begin{frame}{Q44: Risk Reporting Quality}
		\begin{block}{Question}
			Executive risk reports should emphasize:
		\end{block}
		\begin{enumerate}[(A)]
			\item Raw technical data
			\item Business impact, trends, and appetite comparison
			\item Every incident
			\item Vendor pricing
		\end{enumerate}
	\end{frame}
	
	\begin{frame}{Q44: Explanation}
		\begin{block}{Correct Answer: B}
			\textbf{Business impact matters:}
			\begin{itemize}
				\item Executives need decision-relevant data
				\item Trends and appetite comparison drive action
				\item Raw data overwhelms
			\end{itemize}
		\end{block}
		\begin{alertblock}{Critical Trap}
			Business impact drives decision-making.
		\end{alertblock}
	\end{frame}
	
	\begin{frame}{Q45: KRI vs KPI Reporting}
		\begin{block}{Question}
			Which statement about KRIs and KPIs is correct?
		\end{block}
		\begin{enumerate}[(A)]
			\item They are the same
			\item KRIs signal changing exposure; KPIs measure operational performance
			\item KPIs predict risk
			\item KRIs replace KPIs
		\end{enumerate}
	\end{frame}
	
	\begin{frame}{Q45: Explanation}
		\begin{block}{Correct Answer: B}
			\textbf{KRIs signal risk; KPIs measure performance:}
			\begin{itemize}
				\item KRIs are risk-focused
				\item KPIs measure performance toward goals
				\item Both are useful for different purposes
			\end{itemize}
		\end{block}
		\begin{alertblock}{Critical Trap}
			KRI = risk; KPI = performance.
		\end{alertblock}
	\end{frame}
	
	% ============================================================
	\section{5. Risk Culture and Awareness}
	% ============================================================
	
	\begin{frame}{Q46: Culture vs Training}
		\begin{block}{Question}
			The organization runs annual training but executives publicly dismiss risks. Which has greater impact on culture?
		\end{block}
		\begin{enumerate}[(A)]
			\item Training
			\item Executive behavior
			\item Policy documents
			\item Vendor certifications
		\end{enumerate}
	\end{frame}
	
	\begin{frame}{Q46: Explanation}
		\begin{block}{Correct Answer: B}
			\textbf{Behavior beats training:}
			\begin{itemize}
				\item Culture follows leadership behavior
				\item Training alone cannot offset dismissal
				\item Behavior drives decisions
			\end{itemize}
		\end{block}
		\begin{alertblock}{Critical Trap}
			Behavior beats training in shaping culture.
		\end{alertblock}
	\end{frame}
	
	\begin{frame}{Q47: Blame Culture}
		\begin{block}{Question}
			A blame culture exists. What is the MOST likely effect?
		\end{block}
		\begin{enumerate}[(A)]
			\item Fewer risks reported
			\item More honest reporting
			\item Stronger controls
			\item Higher morale
		\end{enumerate}
	\end{frame}
	
	\begin{frame}{Q47: Explanation}
		\begin{block}{Correct Answer: A}
			\textbf{Blame suppresses disclosure:}
			\begin{itemize}
				\item Blame suppresses disclosure
				\item Early warning becomes rare
				\item Silence hides real risks
			\end{itemize}
		\end{block}
		\begin{alertblock}{Critical Trap}
			Blame kills reporting.
		\end{alertblock}
	\end{frame}
	
	\begin{frame}{Q48: Incentive Misalignment}
		\begin{block}{Question}
			Sales bonuses reward revenue without security considerations. What should governance do?
		\end{block}
		\begin{enumerate}[(A)]
			\item Ignore the incentives
			\item Align incentives with risk-aware behavior
			\item Increase sales targets
			\item Eliminate bonuses
		\end{enumerate}
	\end{frame}
	
	\begin{frame}{Q48: Explanation}
		\begin{block}{Correct Answer: B}
			\textbf{Align incentives:}
			\begin{itemize}
				\item Incentives drive behavior
				\item Misalignment creates unmanaged risk
				\item Alignment encourages balanced decisions
			\end{itemize}
		\end{block}
		\begin{alertblock}{Critical Trap}
			Incentives shape risk-taking.
		\end{alertblock}
	\end{frame}
	
	\begin{frame}{Q49: Awareness Program}
		\begin{block}{Question}
			Which is the STRONGEST indicator that a security awareness program is working?
		\end{block}
		\begin{enumerate}[(A)]
			\item Training completion rates are high
			\item Employees spot and report suspicious activity
			\item Training is delivered
			\item Posters are visible
		\end{enumerate}
	\end{frame}
	
	\begin{frame}{Q49: Explanation}
		\begin{block}{Correct Answer: B}
			\textbf{Behavior indicators matter:}
			\begin{itemize}
				\item Behavioral indicators matter
				\item Completion measures activity, not behavior
				\item Reporting shows real awareness
			\end{itemize}
		\end{block}
		\begin{alertblock}{Critical Trap}
			Behavior beats completion rates.
		\end{alertblock}
	\end{frame}
	
	\begin{frame}{Q50: Executive Sponsorship}
		\begin{block}{Question}
			A security initiative lacks executive sponsorship. What is the likely outcome?
		\end{block}
		\begin{enumerate}[(A)]
			\item Success
			\item Reduced engagement and resources
			\item Immediate failure
			\item Better culture
		\end{enumerate}
	\end{frame}
	
	\begin{frame}{Q50: Explanation}
		\begin{block}{Correct Answer: B}
			\textbf{No sponsorship, no traction:}
			\begin{itemize}
				\item Sponsorship drives resources
				\item Unowned initiatives fade
				\item Leadership visibility sustains effort
			\end{itemize}
		\end{block}
		\begin{alertblock}{Critical Trap}
			No sponsorship, no traction.
		\end{alertblock}
	\end{frame}
	
	\begin{frame}{Q51: Communication of Appetite}
		\begin{block}{Question}
			Risk appetite is documented but never communicated to staff. What is the MOST likely result?
		\end{block}
		\begin{enumerate}[(A)]
			\item Alignment improves
			\item Decisions ignore appetite and drift
			\item Staff still follow appetite
			\item Costs fall
		\end{enumerate}
	\end{frame}
	
	\begin{frame}{Q51: Explanation}
		\begin{block}{Correct Answer: B}
			\textbf{Appetite must be known:}
			\begin{itemize}
				\item Appetite must be known to guide choices
				\item Silence produces drift
				\item Communication aligns behavior
			\end{itemize}
		\end{block}
		\begin{alertblock}{Critical Trap}
			Uncommunicated appetite is unmanaged appetite.
		\end{alertblock}
	\end{frame}
	
	\begin{frame}{Q52: Risk Culture Metrics}
		\begin{block}{Question}
			Which metric BEST supports measuring risk culture?
		\end{block}
		\begin{enumerate}[(A)]
			\item Number of training slides
			\item Reported near misses and risk escalations over time
			\item Vendor count
			\item Budget size
		\end{enumerate}
	\end{frame}
	
	\begin{frame}{Q52: Explanation}
		\begin{block}{Correct Answer: B}
			\textbf{Near misses show culture:}
			\begin{itemize}
				\item Near misses reveal culture maturity
				\item Escalation trends indicate honesty
				\item Volume measures activity, not culture
			\end{itemize}
		\end{block}
		\begin{alertblock}{Critical Trap}
			Culture is measured by behavior and reporting.
		\end{alertblock}
	\end{frame}
	
	\begin{frame}{Q53: Leadership Walkaway}
		\begin{block}{Question}
			Leaders ignore agreed risk decisions. What is the MOST likely culture impact?
		\end{block}
		\begin{enumerate}[(A)]
			\item Stronger compliance
			\item Loss of trust and reduced commitment to risk management
			\item Better alignment
			\item Higher training quality
		\end{enumerate}
	\end{frame}
	
	\begin{frame}{Q53: Explanation}
		\begin{block}{Correct Answer: B}
			\textbf{Inconsistency erodes trust:}
			\begin{itemize}
				\item Inconsistency erodes trust
				\item Commitments must be honored
				\item Employees mirror leadership
			\end{itemize}
		\end{block}
		\begin{alertblock}{Critical Trap}
			Inconsistent leadership destroys credibility.
		\end{alertblock}
	\end{frame}
	
	\begin{frame}{Q54: Whistleblower Channels}
		\begin{block}{Question}
			A strong whistleblower channel supports governance by:
		\end{block}
		\begin{enumerate}[(A)]
			\item Discouraging reporting
			\item Providing a safe route to surface wrongdoing
			\item Replacing audits
			\item Eliminating risk
		\end{enumerate}
	\end{frame}
	
	\begin{frame}{Q54: Explanation}
		\begin{block}{Correct Answer: B}
			\textbf{Safe channels surface issues:}
			\begin{itemize}
				\item Safe channels surface issues early
				\item Anonymity reduces fear
				\item Channels complement other controls
			\end{itemize}
		\end{block}
		\begin{alertblock}{Critical Trap}
			Safe channels surface hidden issues.
		\end{alertblock}
	\end{frame}
	
	\begin{frame}{Q55: Cultural Silence}
		\begin{block}{Question}
			Employees believe reporting will not lead to action. What is the MOST likely outcome?
		\end{block}
		\begin{enumerate}[(A)]
			\item Issues are disclosed anyway
			\item Silence, and undetected risks grow
			\item Better training
			\item More auditors
		\end{enumerate}
	\end{frame}
	
	\begin{frame}{Q55: Explanation}
		\begin{block}{Correct Answer: B}
			\textbf{No follow-through, no reporting:}
			\begin{itemize}
				\item Belief that action will not follow discourages reporting
				\item Silence hides exposure
				\item Action builds trust
			\end{itemize}
		\end{block}
		\begin{alertblock}{Critical Trap}
			No follow-through, no reporting.
		\end{alertblock}
	\end{frame}
	
	% ============================================================
	\section{6. Governance of Third Parties}
	% ============================================================
	
	\begin{frame}{Q56: Third-Party Accountability}
		\begin{block}{Question}
			When a service is outsourced, accountability for the data remains with:
		\end{block}
		\begin{enumerate}[(A)]
			\item The vendor
			\item The outsourcing organization
			\item The regulator
			\item The insurance company
		\end{enumerate}
	\end{frame}
	
	\begin{frame}{Q56: Explanation}
		\begin{block}{Correct Answer: B}
			\textbf{Accountability cannot be outsourced:}
			\begin{itemize}
				\item Accountability cannot be outsourced
				\item Vendors can be responsible operationally
				\item The organization remains accountable to regulators
			\end{itemize}
		\end{block}
		\begin{alertblock}{Critical Trap}
			Outsourcing does not transfer accountability.
		\end{alertblock}
	\end{frame}
	
	\begin{frame}{Q57: Vendor Risk Tiering}
		\begin{block}{Question}
			Which is the BEST way to categorize vendor risk?
		\end{block}
		\begin{enumerate}[(A)]
			\item By contract value
			\item By the criticality of services and data access
			\item By the vendor's marketing
			\item By the vendor's age
		\end{enumerate}
	\end{frame}
	
	\begin{frame}{Q57: Explanation}
		\begin{block}{Correct Answer: B}
			\textbf{Criticality drives risk:}
			\begin{itemize}
				\item Criticality drives risk
				\item Data access influences exposure
				\item Contract value does not reflect risk
			\end{itemize}
		\end{block}
		\begin{alertblock}{Critical Trap}
			Tier vendors by criticality and access.
		\end{alertblock}
	\end{frame}
	
	\begin{frame}{Q58: Fourth-Party Risk}
		\begin{block}{Question}
			A critical vendor uses subcontractors. What is this called and why does it matter?
		\end{block}
		\begin{enumerate}[(A)]
			\item First-party risk; it does not matter
			\item Fourth-party risk; visibility is reduced and risk may be higher than expected
			\item Third-party risk; it is the same
			\item No risk
		\end{enumerate}
	\end{frame}
	
	\begin{frame}{Q58: Explanation}
		\begin{block}{Correct Answer: B}
			\textbf{Fourth-party risk adds hidden risk:}
			\begin{itemize}
				\item Fourth-party risk is vendor-of-vendor
				\item It affects reliability and compliance
				\item Visibility is critical
			\end{itemize}
		\end{block}
		\begin{alertblock}{Critical Trap}
			Vendor chains add hidden risk.
		\end{alertblock}
	\end{frame}
	
	\begin{frame}{Q59: Vendor Exit Strategy}
		\begin{block}{Question}
			Why should an exit plan exist before signing a critical vendor contract?
		\end{block}
		\begin{enumerate}[(A)]
			\item To shorten the contract
			\item To reduce dependency and ensure data return and continuity
			\item To save money
			\item To avoid audits
		\end{enumerate}
	\end{frame}
	
	\begin{frame}{Q59: Explanation}
		\begin{block}{Correct Answer: B}
			\textbf{Plan for exit before entering:}
			\begin{itemize}
				\item Exit plans reduce lock-in and recovery risk
				\item Data return is a legal necessity
				\item Continuity depends on early planning
			\end{itemize}
		\end{block}
		\begin{alertblock}{Critical Trap}
			Plan for exit before you enter.
		\end{alertblock}
	\end{frame}
	
	\begin{frame}{Q60: Contractual Controls}
		\begin{block}{Question}
			Which contract clauses MOST support governance of third-party risk?
		\end{block}
		\begin{enumerate}[(A)]
			\item Price and delivery
			\item Right-to-audit, breach notification, and security obligations
			\item Branding and marketing
			\item Coffee vendor preference
		\end{enumerate}
	\end{frame}
	
	\begin{frame}{Q60: Explanation}
		\begin{block}{Correct Answer: B}
			\textbf{Clauses that enable oversight:}
			\begin{itemize}
				\item Audit rights and obligations give oversight
				\item Notification enables timely response
				\item Price does not address risk
			\end{itemize}
		\end{block}
		\begin{alertblock}{Critical Trap}
			Clauses that enable oversight protect you.
		\end{alertblock}
	\end{frame}
	
	\begin{frame}{Q61: Vendor Assurance}
		\begin{block}{Question}
			A vendor provides no independent assurance. What governance response BEST fits?
		\end{block}
		\begin{enumerate}[(A)]
			\item Accept with no additional controls
			\item Require alternative assurance or add compensating controls
			\item Ignore it
			\item Cancel the contract without assessment
		\end{enumerate}
	\end{frame}
	
	\begin{frame}{Q61: Explanation}
		\begin{block}{Correct Answer: B}
			\textbf{Substitute or compensate:}
			\begin{itemize}
				\item Alternative assurance can substitute
				\item Compensating controls reduce exposure
				\item Cancellation needs context
			\end{itemize}
		\end{block}
		\begin{alertblock}{Critical Trap}
			Substitute or compensate, don't ignore.
		\end{alertblock}
	\end{frame}
	
	\begin{frame}{Q62: Vendor Concentration}
		\begin{block}{Question}
			All critical services depend on one provider. What should governance require?
		\end{block}
		\begin{enumerate}[(A)]
			\item More dependency
			\item Concentration risk assessment and contingency plan
			\item No action
			\item A larger contract
		\end{enumerate}
	\end{frame}
	
	\begin{frame}{Q62: Explanation}
		\begin{block}{Correct Answer: B}
			\textbf{Concentration needs strategy:}
			\begin{itemize}
				\item Concentration increases impact
				\item Contingency plans reduce exposure
				\item Contracts alone do not reduce concentration
			\end{itemize}
		\end{block}
		\begin{alertblock}{Critical Trap}
			Concentration needs explicit strategy.
		\end{alertblock}
	\end{frame}
	
	\begin{frame}{Q63: Vendor Governance Cadence}
		\begin{block}{Question}
			Vendor reviews occur only at contract renewal. What is the weakness?
		\end{block}
		\begin{enumerate}[(A)]
			\item Nothing
			\item Reviews are too infrequent to detect drift
			\item Renewal is sufficient
			\item The vendor prefers renewal
		\end{enumerate}
	\end{frame}
	
	\begin{frame}{Q63: Explanation}
		\begin{block}{Correct Answer: B}
			\textbf{Vendors change between renewals:}
			\begin{itemize}
				\item Vendor risk changes over time
				\item Periodic review catches drift
				\item Renewal-only is too late
			\end{itemize}
		\end{block}
		\begin{alertblock}{Critical Trap}
			Vendors change between renewals.
		\end{alertblock}
	\end{frame}
	
	\begin{frame}{Q64: Vendor Incident Notification}
		\begin{block}{Question}
			Which governance clause is MOST valuable for incident response?
		\end{block}
		\begin{enumerate}[(A)]
			\item Timely notification obligations with defined timelines
			\item Free coffee at meetings
			\item Branding rights
			\item Discount clauses
		\end{enumerate}
	\end{frame}
	
	\begin{frame}{Q64: Explanation}
		\begin{block}{Correct Answer: A}
			\textbf{Notice clauses save time:}
			\begin{itemize}
				\item Timely notice enables response
				\item Defined timelines create accountability
				\item Notifications protect both parties
			\end{itemize}
		\end{block}
		\begin{alertblock}{Critical Trap}
			Notice clauses save time and exposure.
		\end{alertblock}
	\end{frame}
	
	\begin{frame}{Q65: Vendor Certification Reliance}
		\begin{block}{Question}
			The vendor holds an ISO 27001 certificate. What should the practitioner do?
		\end{block}
		\begin{enumerate}[(A)]
			\item Stop all oversight
			\item Confirm scope and recency, and continue ongoing oversight
			\item Ignore the certificate
			\item Assume zero risk
		\end{enumerate}
	\end{frame}
	
	\begin{frame}{Q65: Explanation}
		\begin{block}{Correct Answer: B}
			\textbf{Certification is a checkpoint:}
			\begin{itemize}
				\item Certifications show capability, not current performance
				\item Scope and recency matter
				\item Ongoing oversight remains necessary
			\end{itemize}
		\end{block}
		\begin{alertblock}{Critical Trap}
			Certification is a checkpoint, not a guarantee.
		\end{alertblock}
	\end{frame}
	
	% ============================================================
	\section{7. Risk Appetite and Reporting}
	% ============================================================
	
	\begin{frame}{Q66: Appetite vs Tolerance}
		\begin{block}{Question}
			Which distinction is correct?
		\end{block}
		\begin{enumerate}[(A)]
			\item Appetite is tolerated variation; tolerance is intent
			\item Appetite is the level of risk the enterprise is willing to seek or accept; tolerance is acceptable variation
			\item They are the same
			\item Appetite applies only to IT
		\end{enumerate}
	\end{frame}
	
	\begin{frame}{Q66: Explanation}
		\begin{block}{Correct Answer: B}
			\textbf{Appetite vs tolerance:}
			\begin{itemize}
				\item Appetite reflects willingness
				\item Tolerance reflects permitted variance
				\item Both are set by leadership
			\end{itemize}
		\end{block}
		\begin{alertblock}{Critical Trap}
			Appetite = willingness; tolerance = variation.
		\end{alertblock}
	\end{frame}
	
	\begin{frame}{Q67: Board Reporting Purpose}
		\begin{block}{Question}
			What is the MAIN purpose of risk reporting to the board?
		\end{block}
		\begin{enumerate}[(A)]
			\item To show activity
			\item To enable oversight and informed decisions
			\item To compare vendors
			\item To justify the IT budget
		\end{enumerate}
	\end{frame}
	
	\begin{frame}{Q67: Explanation}
		\begin{block}{Correct Answer: B}
			\textbf{Reporting enables oversight:}
			\begin{itemize}
				\item Oversight requires insight
				\item Decisions need context and consequence
				\item Reports enable action
			\end{itemize}
		\end{block}
		\begin{alertblock}{Critical Trap}
			Reporting exists to enable oversight.
		\end{alertblock}
	\end{frame}
	
	\begin{frame}{Q68: KRI Threshold Breach}
		\begin{block}{Question}
			A KRI threshold is breached. What should governance expect?
		\end{block}
		\begin{enumerate}[(A)]
			\item No action
			\item A defined escalation and response process
			\item A new KRI
			\item An audit
		\end{enumerate}
	\end{frame}
	
	\begin{frame}{Q68: Explanation}
		\begin{block}{Correct Answer: B}
			\textbf{Breaches require responses:}
			\begin{itemize}
				\item Thresholds should trigger action
				\item Defined processes avoid confusion
				\item Ignoring signals breaks trust in governance
			\end{itemize}
		\end{block}
		\begin{alertblock}{Critical Trap}
			Breaches require defined responses.
		\end{alertblock}
	\end{frame}
	
	\begin{frame}{Q69: Reporting Aggregation}
		\begin{block}{Question}
			Reporting aggregates all business units into one risk score. What is the RISK of doing this?
		\end{block}
		\begin{enumerate}[(A)]
			\item Simplicity
			\item Losing detail and missing concentrated risk in a specific unit
			\item Reducing effort
			\item Faster reporting
		\end{enumerate}
	\end{frame}
	
	\begin{frame}{Q69: Explanation}
		\begin{block}{Correct Answer: B}
			\textbf{Aggregation can hide concentration:}
			\begin{itemize}
				\item Aggregation can hide concentration
				\item Detail matters for decisions
				\item Summary and detail must coexist
			\end{itemize}
		\end{block}
		\begin{alertblock}{Critical Trap}
			Aggregation can hide concentration.
		\end{alertblock}
	\end{frame}
	
	\begin{frame}{Q70: Loss of Appetite}
		\begin{block}{Question}
			Leaders begin accepting risks well above the stated appetite. What is the MOST likely consequence?
		\end{block}
		\begin{enumerate}[(A)]
			\item Stronger governance
			\item Appetite becomes meaningless and risk-taking becomes uncontrolled
			\item Reduced training
			\item Higher morale
		\end{enumerate}
	\end{frame}
	
	\begin{frame}{Q70: Explanation}
		\begin{block}{Correct Answer: B}
			\textbf{Unenforced appetite is fiction:}
			\begin{itemize}
				\item Appetite must bind behavior
				\item Unbounded acceptance loses discipline
				\item Governance relies on consistent limits
			\end{itemize}
		\end{block}
		\begin{alertblock}{Critical Trap}
			Unenforced appetite is fiction.
		\end{alertblock}
	\end{frame}
	
	\begin{frame}{Q71: Assurance Independence}
		\begin{block}{Question}
			Which report is MOST credible to the board?
		\end{block}
		\begin{enumerate}[(A)]
			\item A self-assessment
			\item An independent assurance report
			\item A vendor's marketing summary
			\item An informal email
		\end{enumerate}
	\end{frame}
	
	\begin{frame}{Q71: Explanation}
		\begin{block}{Correct Answer: B}
			\textbf{Independence supports credibility:}
			\begin{itemize}
				\item Independence supports credibility
				\item Self-assessment can be biased
				\item Marketing is not assurance
			\end{itemize}
		\end{block}
		\begin{alertblock}{Critical Trap}
			Independent assurance carries weight.
		\end{alertblock}
	\end{frame}
	
	\begin{frame}{Q72: Reporting Frequency}
		\begin{block}{Question}
			Which reporting cadence is MOST appropriate for critical risks?
		\end{block}
		\begin{enumerate}[(A)]
			\item Every five years
			\item Defined cadence with escalation for material change
			\item Never
			\item Annually only
		\end{enumerate}
	\end{frame}
	
	\begin{frame}{Q72: Explanation}
		\begin{block}{Correct Answer: B}
			\textbf{Cadence plus escalation:}
			\begin{itemize}
				\item Cadence maintains oversight
				\item Escalation handles urgent changes
				\item Long gaps miss movement
			\end{itemize}
		\end{block}
		\begin{alertblock}{Critical Trap}
			Regular cadence plus escalation.
		\end{alertblock}
	\end{frame}
	
	\begin{frame}{Q73: Audience Tailoring}
		\begin{block}{Question}
			A board report uses the same technical detail as an operations report. What is the MAIN issue?
		\end{block}
		\begin{enumerate}[(A)]
			\item Nothing
			\item Mismatch between audience and detail reduces usefulness
			\item Too many auditors
			\item Too few incidents
		\end{enumerate}
	\end{frame}
	
	\begin{frame}{Q73: Explanation}
		\begin{block}{Correct Answer: B}
			\textbf{Tailor to the audience:}
			\begin{itemize}
				\item Boards need decision-level content
				\item Operations need technical detail
				\item Tailoring improves use
			\end{itemize}
		\end{block}
		\begin{alertblock}{Critical Trap}
			Tailor the message to the audience.
		\end{alertblock}
	\end{frame}
	
	\begin{frame}{Q74: Risk Statement Quality}
		\begin{block}{Question}
			Which risk statement is BEST for reporting?
		\end{block}
		\begin{enumerate}[(A)]
			\item Servers might break
			\item Unpatched customer-facing servers may allow data breaches causing regulatory fines and lost trust
			\item Patch status is 87\%
			\item There are 47 servers
		\end{enumerate}
	\end{frame}
	
	\begin{frame}{Q74: Explanation}
		\begin{block}{Correct Answer: B}
			\textbf{Risk statements link cause, event, consequence:}
			\begin{itemize}
				\item Good risk statements link cause, event, and consequence
				\item Consequences drive decisions
				\item Raw counts do not
			\end{itemize}
		\end{block}
		\begin{alertblock}{Critical Trap}
			Risk statements link cause, event, consequence.
		\end{alertblock}
	\end{frame}
	
	\begin{frame}{Q75: Reporting Bias}
		\begin{block}{Question}
			Reports consistently minimize negative news. What is the MAIN impact?
		\end{block}
		\begin{enumerate}[(A)]
			\item Stronger governance
			\item Loss of trust and poor decisions
			\item Faster remediation
			\item Higher KPIs
		\end{enumerate}
	\end{frame}
	
	\begin{frame}{Q75: Explanation}
		\begin{block}{Correct Answer: B}
			\textbf{Bias undermines oversight:}
			\begin{itemize}
				\item Bias undermines oversight
				\item Leaders cannot act on distorted data
				\item Trust erodes
			\end{itemize}
		\end{block}
		\begin{alertblock}{Critical Trap}
			Minimized news leads to misinformed decisions.
		\end{alertblock}
	\end{frame}
	
	% ============================================================
	\section{8. Business Continuity Governance}
	% ============================================================
	
	\begin{frame}{Q76: BCP Ownership}
		\begin{block}{Question}
			Who should own the business continuity plan?
		\end{block}
		\begin{enumerate}[(A)]
			\item IT operations
			\item Business leaders, with IT support
			\item Vendors
			\item The helpdesk
		\end{enumerate}
	\end{frame}
	
	\begin{frame}{Q76: Explanation}
		\begin{block}{Correct Answer: B}
			\textbf{Business owns continuity:}
			\begin{itemize}
				\item Business owns continuity priorities
				\item IT supports with technology recovery
				\item Ownership drives priorities
			\end{itemize}
		\end{block}
		\begin{alertblock}{Critical Trap}
			Business owns continuity; IT enables it.
		\end{alertblock}
	\end{frame}
	
	\begin{frame}{Q77: BCP Testing}
		\begin{block}{Question}
			A BCP exists but has never been tested. What is governance's BEST action?
		\end{block}
		\begin{enumerate}[(A)]
			\item Assume it works
			\item Require regular testing and reporting of results
			\item Ignore it
			\item Postpone indefinitely
		\end{enumerate}
	\end{frame}
	
	\begin{frame}{Q77: Explanation}
		\begin{block}{Correct Answer: B}
			\textbf{Untested plans are unproven:}
			\begin{itemize}
				\item Untested plans are unreliable
				\item Testing reveals gaps
				\item Reporting supports oversight
			\end{itemize}
		\end{block}
		\begin{alertblock}{Critical Trap}
			Untested plans are unproven.
		\end{alertblock}
	\end{frame}
	
	\begin{frame}{Q78: Continuity Scope}
		\begin{block}{Question}
			A BCP covers only IT systems. What is the MOST likely gap?
		\end{block}
		\begin{enumerate}[(A)]
			\item Too much coverage
			\item Business process continuity, people, and facilities are missing
			\item Too many tests
			\item Overfunded
		\end{enumerate}
	\end{frame}
	
	\begin{frame}{Q78: Explanation}
		\begin{block}{Correct Answer: B}
			\textbf{Continuity is business, not just IT:}
			\begin{itemize}
				\item Continuity is broader than IT
				\item People and facilities matter
				\item Process continuity is the goal
			\end{itemize}
		\end{block}
		\begin{alertblock}{Critical Trap}
			Continuity is business, not just IT.
		\end{alertblock}
	\end{frame}
	
	\begin{frame}{Q79: Crisis Communication}
		\begin{block}{Question}
			Who should lead crisis communications during a major incident?
		\end{block}
		\begin{enumerate}[(A)]
			\item Any engineer
			\item A designated, trained role with governance authority
			\item The vendor
			\item The insurer
		\end{enumerate}
	\end{frame}
	
	\begin{frame}{Q79: Explanation}
		\begin{block}{Correct Answer: B}
			\textbf{Crisis communications needs governance:}
			\begin{itemize}
				\item Communications needs trained ownership
				\item Authority avoids conflicting messages
				\item Vendors and insurers support, not lead
			\end{itemize}
		\end{block}
		\begin{alertblock}{Critical Trap}
			Crisis communications needs governance.
		\end{alertblock}
	\end{frame}
	
	\begin{frame}{Q80: Plan Maintenance}
		\begin{block}{Question}
			The BCP was updated three years ago but operations have changed significantly. What is the governance risk?
		\end{block}
		\begin{enumerate}[(A)]
			\item None
			\item The plan may not reflect the current environment and may fail when needed
			\item The plan is too new
			\item Too much testing
		\end{enumerate}
	\end{frame}
	
	\begin{frame}{Q80: Explanation}
		\begin{block}{Correct Answer: B}
			\textbf{Stale plans fail when needed:}
			\begin{itemize}
				\item Plans must reflect current operations
				\item Changes invalidate assumptions
				\item Regular reviews maintain relevance
			\end{itemize}
		\end{block}
		\begin{alertblock}{Critical Trap}
			Stale plans fail when needed.
		\end{alertblock}
	\end{frame}
	
	% ============================================================
	\section{9. Governance of IT Resources}
	% ============================================================
	
	\begin{frame}{Q81: IT Investment Approval}
		\begin{block}{Question}
			Which principle governs approval of major IT investments?
		\end{block}
		\begin{enumerate}[(A)]
			\item All investments approved by IT only
			\item Investments are approved by appropriate authority and aligned with business objectives
			\item Investments are approved by the vendor
			\item Only by the audit committee
		\end{enumerate}
	\end{frame}
	
	\begin{frame}{Q81: Explanation}
		\begin{block}{Correct Answer: B}
			\textbf{Investments align with objectives:}
			\begin{itemize}
				\item Investments need authority and alignment
				\item Business objectives guide value
				\item Vendor approval is not governance
			\end{itemize}
		\end{block}
		\begin{alertblock}{Critical Trap}
			Investments align with objectives and authority.
		\end{alertblock}
	\end{frame}
	
	\begin{frame}{Q82: Resource Optimization}
		\begin{block}{Question}
			Resource optimization in IT governance means:
		\end{block}
		\begin{enumerate}[(A)]
			\item Reducing cost at all costs
			\item Balancing cost, risk, and value to meet business objectives
			\item Maximizing headcount
			\item Buying the newest tools
		\end{enumerate}
	\end{frame}
	
	\begin{frame}{Q82: Explanation}
		\begin{block}{Correct Answer: B}
			\textbf{Balance cost, risk, and value:}
			\begin{itemize}
				\item Optimization is balanced
				\item Value and risk matter alongside cost
				\item Cost reduction alone creates risk
			\end{itemize}
		\end{block}
		\begin{alertblock}{Critical Trap}
			Balance cost, risk, and value.
		\end{alertblock}
	\end{frame}
	
	\begin{frame}{Q83: Performance Measurement}
		\begin{block}{Question}
			IT performance should be measured against:
		\end{block}
		\begin{enumerate}[(A)]
			\item Peer opinions
			\item Business objectives and agreed service levels
			\item Vendor promises
			\item Employee satisfaction alone
		\end{enumerate}
	\end{frame}
	
	\begin{frame}{Q83: Explanation}
		\begin{block}{Correct Answer: B}
			\textbf{Measure against business goals:}
			\begin{itemize}
				\item Performance aligns with business goals
				\item Service levels define expectations
				\item Opinions and promises are not measures
			\end{itemize}
		\end{block}
		\begin{alertblock}{Critical Trap}
			Measure IT against business goals.
		\end{alertblock}
	\end{frame}
	
	\begin{frame}{Q84: Value Delivery}
		\begin{block}{Question}
			Value delivery from IT is demonstrated when:
		\end{block}
		\begin{enumerate}[(A)]
			\item Projects finish on budget alone
			\item Benefits are realized and aligned with strategy
			\item Technology is modern
			\item Licenses are cheap
		\end{enumerate}
	\end{frame}
	
	\begin{frame}{Q84: Explanation}
		\begin{block}{Correct Answer: B}
			\textbf{Value = realized benefits:}
			\begin{itemize}
				\item Value is realized benefits
				\item Strategy alignment matters
				\item Budget and modernity do not equal value
			\end{itemize}
		\end{block}
		\begin{alertblock}{Critical Trap}
			Value = realized benefits aligned with strategy.
		\end{alertblock}
	\end{frame}
	
	\begin{frame}{Q85: Governance of Change}
		\begin{block}{Question}
			Which control best supports governance over significant IT changes?
		\end{block}
		\begin{enumerate}[(A)]
			\item Every change is approved
			\item A defined change process with appropriate authority levels
			\item Only emergency changes are governed
			\item Vendors approve all changes
		\end{enumerate}
	\end{frame}
	
	\begin{frame}{Q85: Explanation}
		\begin{block}{Correct Answer: B}
			\textbf{Authority-based change control:}
			\begin{itemize}
				\item Authority-based change control fits risk
				\item Excessive approvals slow operations
				\item Vendors must not approve internal changes
			\end{itemize}
		\end{block}
		\begin{alertblock}{Critical Trap}
			Change governance is authority-based.
		\end{alertblock}
	\end{frame}
	
	% ============================================================
	\section{10. Advanced Governance Scenarios}
	% ============================================================
	
	\begin{frame}{Q86: Board Knowledge Gap}
		\begin{block}{Question}
			The board has no IT expertise. What is the BEST governance response?
		\end{block}
		\begin{enumerate}[(A)]
			\item Ignore IT risk
			\item Establish training, independent advice, or a technology committee to support oversight
			\item Delegate all IT to the CIO
			\item Cancel IT projects
		\end{enumerate}
	\end{frame}
	
	\begin{frame}{Q86: Explanation}
		\begin{block}{Correct Answer: B}
			\textbf{Capability for oversight:}
			\begin{itemize}
				\item Oversight requires capability
				\item Advisory structures support boards
				\item Deferring IT entirely abdicates oversight
			\end{itemize}
		\end{block}
		\begin{alertblock}{Critical Trap}
			Boards need capability to oversee IT.
		\end{alertblock}
	\end{frame}
	
	\begin{frame}{Q87: Board Independence}
		\begin{block}{Question}
			The board majority is composed of executives from the company. What is the concern?
		\end{block}
		\begin{enumerate}[(A)]
			\item Stronger oversight
			\item Weakened independence and potential conflict of interest
			\item Faster decisions
			\item Lower cost
		\end{enumerate}
	\end{frame}
	
	\begin{frame}{Q87: Explanation}
		\begin{block}{Correct Answer: B}
			\textbf{Independence protects oversight:}
			\begin{itemize}
				\item Independence supports oversight
				\item Executive-heavy boards may lack challenge
				\item Governance depends on independence
			\end{itemize}
		\end{block}
		\begin{alertblock}{Critical Trap}
			Independence protects oversight.
		\end{alertblock}
	\end{frame}
	
	\begin{frame}{Q88: Executive Removal}
		\begin{block}{Question}
			An executive was dismissed for reporting a serious risk. What is the MOST likely culture outcome?
		\end{block}
		\begin{enumerate}[(A)]
			\item Stronger risk management
			\item Silence and suppressed risk reporting
			\item Better training
			\item Increased audits
		\end{enumerate}
	\end{frame}
	
	\begin{frame}{Q88: Explanation}
		\begin{block}{Correct Answer: B}
			\textbf{Punishment silences:}
			\begin{itemize}
				\item Fear drives silence
				\item Retaliation damages trust
				\item Culture erodes
			\end{itemize}
		\end{block}
		\begin{alertblock}{Critical Trap}
			Punishing risk reporting silences the organization.
		\end{alertblock}
	\end{frame}
	
	\begin{frame}{Q89: Governance of Outages}
		\begin{block}{Question}
			Repeated outages are tolerated without change. What governance principle is MOST relevant?
		\end{block}
		\begin{enumerate}[(A)]
			\item Ignoring trends
			\item Accountability for performance and corrective action
			\item Increasing budgets
			\item Marketing better
		\end{enumerate}
	\end{frame}
	
	\begin{frame}{Q89: Explanation}
		\begin{block}{Correct Answer: B}
			\textbf{Repeated failures need response:}
			\begin{itemize}
				\item Governance requires performance ownership
				\item Repeated failures signal gaps
				\item Corrective action is expected
			\end{itemize}
		\end{block}
		\begin{alertblock}{Critical Trap}
			Repeated failures need governance response.
		\end{alertblock}
	\end{frame}
	
	\begin{frame}{Q90: Mandate Clarity}
		\begin{block}{Question}
			A regulatory mandate conflicts with a business objective. What is the BEST governance action?
		\end{block}
		\begin{enumerate}[(A)]
			\item Ignore the regulation
			\item Identify the conflict and escalate to leadership for resolution
			\item Follow the business
			\item Delay indefinitely
		\end{enumerate}
	\end{frame}
	
	\begin{frame}{Q90: Explanation}
		\begin{block}{Correct Answer: B}
			\textbf{Conflicts go to leadership:}
			\begin{itemize}
				\item Conflicts require leadership decisions
				\item Legal obligations usually constrain business choices
				\item Escalation is the right path
			\end{itemize}
		\end{block}
		\begin{alertblock}{Critical Trap}
			Conflicts go to leadership for resolution.
		\end{alertblock}
	\end{frame}
	
	\begin{frame}{Q91: Reporting Lines}
		\begin{block}{Question}
			The risk function reports directly to the CRO who reports to the CEO. What is the governance advantage?
		\end{block}
		\begin{enumerate}[(A)]
			\item Less independence
			\item Stronger independence and visibility of risk
			\item Lower cost
			\item Faster operations
		\end{enumerate}
	\end{frame}
	
	\begin{frame}{Q91: Explanation}
		\begin{block}{Correct Answer: B}
			\textbf{Direct lines strengthen independence:}
			\begin{itemize}
				\item Direct lines improve visibility
				\item Independence supports honesty
				\item Cost and speed are not the point
			\end{itemize}
		\end{block}
		\begin{alertblock}{Critical Trap}
			Direct lines strengthen independence.
		\end{alertblock}
	\end{frame}
	
	\begin{frame}{Q92: Committee Composition}
		\begin{block}{Question}
			A risk committee lacks legal and finance members. What governance risk is MOST likely?
		\end{block}
		\begin{enumerate}[(A)]
			\item Faster decisions
			\item Blind spots in legal and financial exposure
			\item Better technology focus
			\item Lower costs
		\end{enumerate}
	\end{frame}
	
	\begin{frame}{Q92: Explanation}
		\begin{block}{Correct Answer: B}
			\textbf{Missing voices mean missing risks:}
			\begin{itemize}
				\item Cross-functional perspectives catch blind spots
				\item Legal and finance affect risk materiality
				\item One-sided committees under-represent risk
			\end{itemize}
		\end{block}
		\begin{alertblock}{Critical Trap}
			Missing voices mean missing risks.
		\end{alertblock}
	\end{frame}
	
	\begin{frame}{Q93: Governance of Emerging Tech}
		\begin{block}{Question}
			A new technology is rapidly adopted across business units. What governance is MOST needed?
		\end{block}
		\begin{enumerate}[(A)]
			\item Slower adoption
			\item A defined framework for assessment, ownership, and oversight
			\item Immediate shutdown
			\item No governance
		\end{enumerate}
	\end{frame}
	
	\begin{frame}{Q93: Explanation}
		\begin{block}{Correct Answer: B}
			\textbf{Adoption needs a framework:}
			\begin{itemize}
				\item Frameworks balance adoption and risk
				\item Ownership ensures accountability
				\item Shutdown may harm the business
			\end{itemize}
		\end{block}
		\begin{alertblock}{Critical Trap}
			Adoption needs a governance framework.
		\end{alertblock}
	\end{frame}
	
	\begin{frame}{Q94: Decision Rights}
		\begin{block}{Question}
			Which document best clarifies who makes decisions?
		\end{block}
		\begin{enumerate}[(A)]
			\item The vendor contract
			\item A RACI chart or decision rights matrix
			\item The audit report
			\item A marketing plan
		\end{enumerate}
	\end{frame}
	
	\begin{frame}{Q94: Explanation}
		\begin{block}{Correct Answer: B}
			\textbf{Decision rights belong in RACI:}
			\begin{itemize}
				\item Decision rights clarify authority
				\item RACI assigns roles
				\item Other documents do not assign decision authority
			\end{itemize}
		\end{block}
		\begin{alertblock}{Critical Trap}
			Decision rights belong in RACI or matrices.
		\end{alertblock}
	\end{frame}
	
	\begin{frame}{Q95: Performance Escalation}
		\begin{block}{Question}
			Business unit performance repeatedly misses risk commitments. What is the BEST governance action?
		\end{block}
		\begin{enumerate}[(A)]
			\item Ignore the misses
			\item Escalate and require corrective action with accountability
			\item Raise budgets
			\item Cancel training
		\end{enumerate}
	\end{frame}
	
	\begin{frame}{Q95: Explanation}
		\begin{block}{Correct Answer: B}
			\textbf{Missed commitments need escalation:}
			\begin{itemize}
				\item Escalation addresses accountability
				\item Corrective action follows
				\item Ignoring normalizes failure
			\end{itemize}
		\end{block}
		\begin{alertblock}{Critical Trap}
			Missed commitments need escalation.
		\end{alertblock}
	\end{frame}
	
	\begin{frame}{Q96: Board Risk Reporting}
		\begin{block}{Question}
			Which content is MOST valuable in a board risk report?
		\end{block}
		\begin{enumerate}[(A)]
			\item Detailed technical findings
			\item Key exposures, appetite comparison, and recommended actions
			\item Vendor contract details
			\item Marketing performance
		\end{enumerate}
	\end{frame}
	
	\begin{frame}{Q96: Explanation}
		\begin{block}{Correct Answer: B}
			\textbf{Boards need exposures and actions:}
			\begin{itemize}
				\item Boards need exposures and recommendations
				\item Appetite comparison guides oversight
				\item Technical and commercial details overwhelm
			\end{itemize}
		\end{block}
		\begin{alertblock}{Critical Trap}
			Boards need exposures and actions.
		\end{alertblock}
	\end{frame}
	
	\begin{frame}{Q97: Risk Acceptance Formality}
		\begin{block}{Question}
			A business owner accepts a large risk without formal documentation. What governance gap is MOST likely?
		\end{block}
		\begin{enumerate}[(A)]
			\item Missing documentation and accountability
			\item Too many documents
			\item Too many auditors
			\item Too little technology
		\end{enumerate}
	\end{frame}
	
	\begin{frame}{Q97: Explanation}
		\begin{block}{Correct Answer: A}
			\textbf{Formal acceptance is governance:}
			\begin{itemize}
				\item Documentation creates accountability
				\item Informal acceptance is indefensible
				\item Governance requires traceability
			\end{itemize}
		\end{block}
		\begin{alertblock}{Critical Trap}
			Formal acceptance is governance.
		\end{alertblock}
	\end{frame}
	
	\begin{frame}{Q98: Conflict of Interest}
		\begin{block}{Question}
			A risk committee member owns shares in a vendor under review. What should happen?
		\end{block}
		\begin{enumerate}[(A)]
			\item Proceed without disclosure
			\item Disclose and recuse from related decisions
			\item Vote on the vendor
			\item Ignore the conflict
		\end{enumerate}
	\end{frame}
	
	\begin{frame}{Q98: Explanation}
		\begin{block}{Correct Answer: B}
			\textbf{Disclose and recuse:}
			\begin{itemize}
				\item Conflicts must be disclosed
				\item Recusal protects integrity
				\item Silence damages trust
			\end{itemize}
		\end{block}
		\begin{alertblock}{Critical Trap}
			Disclose conflicts and recuse.
		\end{alertblock}
	\end{frame}
	
	\begin{frame}{Q99: Executive Reporting}
		\begin{block}{Question}
			An executive proposes reducing the risk budget without analysis. What is the BEST governance response?
		\end{block}
		\begin{enumerate}[(A)]
			\item Approve immediately
			\item Require risk impact analysis and a formal decision
			\item Delay indefinitely
			\item Ignore the proposal
		\end{enumerate}
	\end{frame}
	
	\begin{frame}{Q99: Explanation}
		\begin{block}{Correct Answer: B}
			\textbf{Budget cuts need risk analysis:}
			\begin{itemize}
				\item Budget cuts affect risk exposure
				\item Analysis supports informed decisions
				\item Formal decision records accountability
			\end{itemize}
		\end{block}
		\begin{alertblock}{Critical Trap}
			Budget cuts need risk analysis.
		\end{alertblock}
	\end{frame}
	
	\begin{frame}{Q100: Ultimate Accountability}
		\begin{block}{Question}
			Ultimate accountability for enterprise risk rests with:
		\end{block}
		\begin{enumerate}[(A)]
			\item The risk practitioner
			\item The CISO
			\item The board of directors
			\item The vendor
		\end{enumerate}
	\end{frame}
	
	\begin{frame}{Q100: Explanation}
		\begin{block}{Correct Answer: C}
			\textbf{Ultimate accountability sits with the board:}
			\begin{itemize}
				\item Boards are ultimately accountable to stakeholders
				\item Management executes within direction
				\item Practitioners and CISOs advise and operate
			\end{itemize}
		\end{block}
		\begin{alertblock}{Critical Trap}
			Ultimate accountability sits with the board.
		\end{alertblock}
	\end{frame}
	
	\begin{frame}{Exam Mindset Summary}
		\begin{block}{Remember}
			\begin{itemize}
				\item \textbf{Governance} directs; \textbf{management} executes
				\item \textbf{Accountability} cannot be delegated or outsourced
				\item \textbf{Appetite} is board-owned; \textbf{tolerance} is permitted variation
				\item \textbf{Three lines:} operate (1st), oversee (2nd), audit (3rd)
				\item \textbf{Policy} states intent; \textbf{standards} specify requirements
				\item \textbf{Culture} follows leadership behavior, not slogans
			\end{itemize}
		\end{block}
	\end{frame}
	
	\begin{frame}{Thank You}
		\begin{center}
			{\Huge \textbf{Good Luck on Your CRISC Exam!}}
		\end{center}
	\end{frame}
	
\end{document}